% =====================================================================
%  probe.tex — renders every figure in figures/src/ on its own page, so
%  a human can look at all six before they go anywhere near a chapter.
%
%  Run from the thesis root, so that \input{style/...} resolves:
%
%      xelatex -output-directory=figures/out figures/probe.tex
%      xelatex -output-directory=figures/out "\def\probelang{el}% =====================================================================
%  probe.tex — renders every figure in figures/src/ on its own page, so
%  a human can look at all six before they go anywhere near a chapter.
%
%  Run from the thesis root, so that \input{style/...} resolves:
%
%      xelatex -output-directory=figures/out figures/probe.tex
%      xelatex -output-directory=figures/out "\def\probelang{el}% =====================================================================
%  probe.tex — renders every figure in figures/src/ on its own page, so
%  a human can look at all six before they go anywhere near a chapter.
%
%  Run from the thesis root, so that \input{style/...} resolves:
%
%      xelatex -output-directory=figures/out figures/probe.tex
%      xelatex -output-directory=figures/out "\def\probelang{el}% =====================================================================
%  probe.tex — renders every figure in figures/src/ on its own page, so
%  a human can look at all six before they go anywhere near a chapter.
%
%  Run from the thesis root, so that \input{style/...} resolves:
%
%      xelatex -output-directory=figures/out figures/probe.tex
%      xelatex -output-directory=figures/out "\def\probelang{el}\input{figures/probe}"
%
%  Both editions must compile with zero errors AND be looked at:
%      pdftoppm -png -r 110 figures/out/probe.pdf /tmp/probe
% =====================================================================
\documentclass[11pt,a4paper]{report}
\providecommand{\probelang}{en}
\input{style/common}
\input{style/lang-\probelang}
\input{figures/labels-\probelang}

\pagestyle{empty}
\begin{document}
\thesisnames

\newcommand{\probe}[1]{%
  \begin{center}
    {\ttfamily\small figures/src/#1.tex \quad(\probelang)}\par\vspace{6mm}
    \input{figures/src/#1}
  \end{center}
  \clearpage}

\probe{fig-timeline}
\probe{fig-quorum}
\probe{fig-scaling}
\probe{fig-synchrony}
\probe{fig-equivocation}
\probe{fig-smr}
\end{document}
"
%
%  Both editions must compile with zero errors AND be looked at:
%      pdftoppm -png -r 110 figures/out/probe.pdf /tmp/probe
% =====================================================================
\documentclass[11pt,a4paper]{report}
\providecommand{\probelang}{en}
% =====================================================================
%  common.tex — everything that is the same in both languages.
%
%  Compile with XeLaTeX. pdfLaTeX will not do: the Greek edition needs
%  fontspec, and both editions must use the same font files so the two
%  PDFs look like the same document.
%
%  Language-specific setup lives in style/lang-en.tex and lang-el.tex.
% =====================================================================

% ---- Fonts ----------------------------------------------------------
\usepackage{fontspec}
\usepackage{polyglossia}

% The department requires Times New Roman at 11 pt (Guidelines A.2).
% Times New Roman itself is not redistributable, so the build asks for it
% first and falls back to Liberation Serif, which is metric-compatible:
% same widths, same line breaks, same page count. On a machine that has
% the real font (Word, Overleaf) the first branch wins and nothing moves.
\IfFontExistsTF{Times New Roman}{
  \newcommand{\thesisserif}{Times New Roman}
  \newcommand{\thesissans}{Arial}
  \newcommand{\thesismono}{Courier New}
}{
  \newcommand{\thesisserif}{Liberation Serif}
  \newcommand{\thesissans}{Liberation Sans}
  \newcommand{\thesismono}{Liberation Mono}
}
\setmainfont[Ligatures=TeX]{\thesisserif}
\setsansfont[Ligatures=TeX]{\thesissans}
\setmonofont[Scale=0.90,Ligatures=TeX]{\thesismono}

% ---- Page ------------------------------------------------------------
% 2.54 cm on every side — the margin the department template uses.
\usepackage[a4paper,margin=2.54cm]{geometry}
\usepackage{microtype}
\emergencystretch=1.5em          % prefer loose spacing over an overfull line
\usepackage{setspace}
\onehalfspacing                  % Guidelines A.2: line spacing 1.5
\setlength{\parskip}{0pt}
\setlength{\parindent}{1.2em}

\usepackage{csquotes}
\usepackage{enumitem}
\setlist{itemsep=2pt,parsep=0pt,topsep=4pt}

% ---- Tables and figures ---------------------------------------------
\usepackage{booktabs}
\usepackage{tabularx}
\usepackage{array}
\usepackage{multirow}
\usepackage{graphicx}
\usepackage{xcolor}
\usepackage[font=small,labelfont=bf,textfont=small]{caption}
\usepackage{subcaption}
\usepackage{float}

% ---- Diagrams --------------------------------------------------------
\usepackage{tikz}
\usetikzlibrary{positioning,arrows.meta,calc,fit,backgrounds,shapes.geometric,
                decorations.pathreplacing}
\usepackage{pgfplots}
\pgfplotsset{compat=1.18}
\usepgfplotslibrary{groupplots}
\input{style/diagrams}

% ---- Maths -----------------------------------------------------------
\usepackage{amsmath,amssymb}

% ---- Citations and links ---------------------------------------------
% cite sorts and compresses numeric citations: \cite{a,b,c} -> [2]-[4].
% It must load BEFORE hyperref.
\usepackage{cite}
\usepackage[colorlinks=true,linkcolor=black,citecolor=black,
            urlcolor=dgCFT,breaklinks=true]{hyperref}
\usepackage{url}
\urlstyle{same}

% ---- Source-quality markers ------------------------------------------
% Project rule R9: a source that has not been through peer review is
% marked WHERE IT IS USED, not only in the bibliography. A reader
% checking a claim should not have to turn to the back to find out that
% the evidence is a preprint.
% Capitals, not small capitals: Times New Roman has no true small-caps
% face, and neither does the metric-compatible fallback, so \textsc{}
% only produced a font warning and a synthesised shape.
\newcommand{\grey}{\textsuperscript{G}}   % grey literature
\newcommand{\prep}{\textsuperscript{P}}   % preprint, not refereed
% U marks a DIFFERENT thing from G and P. Those two say what kind of
% publication a source is. U says that the thesis could not verify the
% source's own bibliographic data against the article or an index, and
% took it from another work's reference list. A reader who wants to
% follow the citation needs to know that before they try.
\newcommand{\unver}{\textsuperscript{U}}   % bibliographic data unverified

% ---- Front-matter heading -------------------------------------------
% 14 pt bold, centred — the department template's look for the pages
% before Chapter 1. \chapter* would print a three-line monster instead.
%   #1 = as printed   #2 = as it appears in the table of contents
\newcommand{\frontheading}[2]{%
  \phantomsection
  \addcontentsline{toc}{chapter}{#2}%
  \markboth{#2}{#2}%
  \begin{center}\large\bfseries #1\end{center}
  \vspace{0.8em}}

% Blank leaf, as the template has between front-matter sections. It is
% counted and numbered, so both editions and the .docx agree page for page.
\newcommand{\blankleaf}{%
  \clearpage
  \thispagestyle{plain}%
  \null\vfill\begin{center}\small\dgblankleaftext\end{center}\vfill\null
  \clearpage}

% ---- Measurements that do not exist yet ------------------------------
% Project rule R1: never print a number that was not measured. Every
% figure from the author's own benchmark campaign is declared once in
% data/measurements.tex. Until a value is filled in, \meas prints a
% visible marker instead of a plausible-looking number, and
% tools/data-check.py lists what is still outstanding.
\newcommand{\PENDING}{\textbf{\color{dgBFT}[PENDING]}}
\newcommand{\meas}[1]{\csname measvalue@#1\endcsname}
\newcommand{\declaremeasurement}[2]{%
  \expandafter\newcommand\csname measvalue@#1\endcsname{#2}}
\newcommand{\pendingmeasurement}[1]{%
  \expandafter\newcommand\csname measvalue@#1\endcsname{\PENDING}}

\input{style/lang-\probelang}
\input{figures/labels-\probelang}

\pagestyle{empty}
\begin{document}
\thesisnames

\newcommand{\probe}[1]{%
  \begin{center}
    {\ttfamily\small figures/src/#1.tex \quad(\probelang)}\par\vspace{6mm}
    \input{figures/src/#1}
  \end{center}
  \clearpage}

\probe{fig-timeline}
\probe{fig-quorum}
\probe{fig-scaling}
\probe{fig-synchrony}
\probe{fig-equivocation}
\probe{fig-smr}
\end{document}
"
%
%  Both editions must compile with zero errors AND be looked at:
%      pdftoppm -png -r 110 figures/out/probe.pdf /tmp/probe
% =====================================================================
\documentclass[11pt,a4paper]{report}
\providecommand{\probelang}{en}
% =====================================================================
%  common.tex — everything that is the same in both languages.
%
%  Compile with XeLaTeX. pdfLaTeX will not do: the Greek edition needs
%  fontspec, and both editions must use the same font files so the two
%  PDFs look like the same document.
%
%  Language-specific setup lives in style/lang-en.tex and lang-el.tex.
% =====================================================================

% ---- Fonts ----------------------------------------------------------
\usepackage{fontspec}
\usepackage{polyglossia}

% The department requires Times New Roman at 11 pt (Guidelines A.2).
% Times New Roman itself is not redistributable, so the build asks for it
% first and falls back to Liberation Serif, which is metric-compatible:
% same widths, same line breaks, same page count. On a machine that has
% the real font (Word, Overleaf) the first branch wins and nothing moves.
\IfFontExistsTF{Times New Roman}{
  \newcommand{\thesisserif}{Times New Roman}
  \newcommand{\thesissans}{Arial}
  \newcommand{\thesismono}{Courier New}
}{
  \newcommand{\thesisserif}{Liberation Serif}
  \newcommand{\thesissans}{Liberation Sans}
  \newcommand{\thesismono}{Liberation Mono}
}
\setmainfont[Ligatures=TeX]{\thesisserif}
\setsansfont[Ligatures=TeX]{\thesissans}
\setmonofont[Scale=0.90,Ligatures=TeX]{\thesismono}

% ---- Page ------------------------------------------------------------
% 2.54 cm on every side — the margin the department template uses.
\usepackage[a4paper,margin=2.54cm]{geometry}
\usepackage{microtype}
\emergencystretch=1.5em          % prefer loose spacing over an overfull line
\usepackage{setspace}
\onehalfspacing                  % Guidelines A.2: line spacing 1.5
\setlength{\parskip}{0pt}
\setlength{\parindent}{1.2em}

\usepackage{csquotes}
\usepackage{enumitem}
\setlist{itemsep=2pt,parsep=0pt,topsep=4pt}

% ---- Tables and figures ---------------------------------------------
\usepackage{booktabs}
\usepackage{tabularx}
\usepackage{array}
\usepackage{multirow}
\usepackage{graphicx}
\usepackage{xcolor}
\usepackage[font=small,labelfont=bf,textfont=small]{caption}
\usepackage{subcaption}
\usepackage{float}

% ---- Diagrams --------------------------------------------------------
\usepackage{tikz}
\usetikzlibrary{positioning,arrows.meta,calc,fit,backgrounds,shapes.geometric,
                decorations.pathreplacing}
\usepackage{pgfplots}
\pgfplotsset{compat=1.18}
\usepgfplotslibrary{groupplots}
% =====================================================================
%  diagrams.tex — THE diagram standard. One file, every figure.
%
%  The rule this file exists to enforce: a reader should be able to look
%  at any two figures in the thesis and see the same visual language.
%  Same palette, same line weights, same type size, same width.
%
%  There are exactly TWO kinds of figure, and every figure is one of them:
%
%    dgdata  — numbers on axes (pgfplots). Anything measured or computed.
%    dgflow  — structure and sequence (TikZ). Anything explained.
%
%  Nothing in content/ may define a colour, a width, or a font size.
%  tools/diagram-check.py fails the build if it does.
% =====================================================================

% ---------------------------------------------------------------------
%  Palette — Okabe & Ito, chosen because it survives the three common
%  colour-vision deficiencies AND greyscale printing. A thesis gets
%  printed in black and white more often than anyone admits.
%
%  Each colour has ONE meaning across the whole thesis:
% ---------------------------------------------------------------------
\definecolor{dgCFT}{HTML}{0072B2}   % crash-fault protocols            (blue)
\definecolor{dgBFT}{HTML}{D55E00}   % Byzantine-fault protocols     (vermilion)
\definecolor{dgAux}{HTML}{009E73}   % the third series / "correct"     (green)
\definecolor{dgSand}{HTML}{E69F00}  % the fourth series / "attention" (orange)
\definecolor{dgPurple}{HTML}{CC79A7}% the fifth series                (purple)
\definecolor{dgGrey}{HTML}{6E6E6E}  % structure only — never a series   (grey)
\definecolor{dgInk}{HTML}{1A1A1A}   % text and axes
\definecolor{dgPaper}{HTML}{FFFFFF} % label backgrounds over lines

% Colour is never the only difference between two series. Every series
% also gets its own mark and its own dash pattern, so the figure still
% reads on a photocopy. These four cycles stay in lockstep.
\pgfplotscreateplotcyclelist{dgcycle}{%
  {dgCFT,   mark=*,        mark size=1.9pt, solid},
  {dgBFT,   mark=square*,  mark size=1.8pt, dashed},
  {dgAux,   mark=triangle*,mark size=2.1pt, dotted},
  {dgSand,  mark=diamond*, mark size=2.1pt, dashdotted},
  {dgPurple,mark=pentagon*,mark size=1.9pt, densely dashed},
}

% ---------------------------------------------------------------------
%  Fixed geometry. Every figure is the same width, so the eye does not
%  have to re-scale between them.
% ---------------------------------------------------------------------
% Computed at \begin{document}, NOT here: in the preamble \linewidth is
% still the class default (360pt), not the real text block set by
% geometry. Setting it here silently made every figure 10.9cm wide on a
% 15.9cm page.
\newlength{\dgwidth}
\newlength{\dgheight}
\AtBeginDocument{%
  \setlength{\dgwidth}{0.86\linewidth}%
  \setlength{\dgheight}{0.50\dgwidth}%
}

% ---------------------------------------------------------------------
%  KIND 1 — dgdata: numbers on axes
%
%    \begin{dgdata}{xlabel={...}, ylabel={...}}
%      \addplot ... ;
%    \end{dgdata}
%
%  Everything that could drift between figures is set here once.
% ---------------------------------------------------------------------
\pgfplotsset{
  dgbase/.style={
    width=\dgwidth, height=\dgheight,
    cycle list name=dgcycle,
    every axis plot/.append style={line width=0.9pt},
    tick label style={font=\footnotesize, color=dgInk},
    label style={font=\small, color=dgInk},
    title style={font=\small\bfseries, color=dgInk},
    legend style={font=\footnotesize, draw=dgGrey!45, fill=dgPaper,
                  fill opacity=0.92, text opacity=1,
                  cells={anchor=west}, row sep=1pt},
    legend cell align=left,
    axis line style={color=dgGrey!70, line width=0.6pt},
    grid=major,
    major grid style={color=dgGrey!22, line width=0.4pt},
    tick align=outside, tick style={color=dgGrey!70},
    enlarge x limits=0.06, enlarge y limits=0.10,
    axis on top,
  },
  % Bar charts share the same skin but need their own spacing.
  dgbars/.style={
    dgbase,
    ybar, bar width=8pt, enlarge x limits=0.14,
    ymin=0, grid=major, ymajorgrids=true, xmajorgrids=false,
    % "fill=." means "fill with this plot's own cycle colour". Without it
    % the cycle list sets only the line colour and every bar comes out as
    % an empty outline.
    every axis plot/.append style={fill=., fill opacity=0.85,
                                   draw=dgInk!55, line width=0.4pt},
    xtick style={draw=none},
  },
}

\newenvironment{dgdata}[1][]
  {\begin{tikzpicture}\begin{axis}[dgbase,#1]}
  {\end{axis}\end{tikzpicture}}

\newenvironment{dgbarchart}[1][]
  {\begin{tikzpicture}\begin{axis}[dgbars,#1]}
  {\end{axis}\end{tikzpicture}}

% ---------------------------------------------------------------------
%  KIND 2 — dgflow: structure and sequence
%
%    \begin{dgflow}
%      \node[dgbox] (a) {...};
%      \node[dgbox, right=of a] (b) {...};
%      \draw[dgarrow] (a) -- (b) node[dgtag,midway]{...};
%    \end{dgflow}
%
%  Node styles are deliberately few. Five box styles is already too many
%  for a reader to hold; these are four, and each means one thing.
% ---------------------------------------------------------------------
\tikzset{
  dgflowbase/.style={
    node distance=9mm and 13mm,
    every node/.append style={font=\small, text=dgInk, align=center},
    >={Stealth[length=2.4mm,width=1.8mm]},
  },
  % a thing that exists: a node, a component, a layer
  dgbox/.style={
    draw=dgGrey!75, line width=0.6pt, rounded corners=1.5pt,
    fill=dgGrey!7, inner sep=4pt, minimum height=7mm, minimum width=20mm,
  },
  % the same, but this one is the subject of the figure
  dgboxkey/.style={dgbox, draw=dgCFT, line width=0.9pt, fill=dgCFT!8},
  % the same, but this one is where it goes wrong
  dgboxbad/.style={dgbox, draw=dgBFT, line width=0.9pt, fill=dgBFT!8},
  % a moment in time, or a state
  dgstate/.style={
    draw=dgGrey!75, line width=0.6pt, circle, fill=dgPaper,
    inner sep=2.5pt, minimum size=9mm,
  },
  % movement: a message, a step, a dependency
  dgarrow/.style={->, draw=dgGrey!85, line width=0.7pt},
  dgarrowkey/.style={dgarrow, draw=dgCFT, line width=0.9pt},
  dgarrowbad/.style={dgarrow, draw=dgBFT, line width=0.9pt, dashed},
  % a label sitting on a line — the white background is what stops it
  % colliding with the line underneath
  % NOTE: text=, not color=. TikZ's color= sets the fill too, so
  % "color=dgInk" here would override fill=dgPaper and render every tag
  % as dark text on a dark box.
  dgtag/.style={
    font=\footnotesize, text=dgInk, fill=dgPaper, inner sep=1.5pt,
  },
  % an aside: the sentence that explains what the reader is looking at
  dgnote/.style={
    font=\footnotesize\itshape, color=dgGrey, align=left, inner sep=2pt,
  },
  % an axis inside a flow figure (a timeline, a spectrum)
  dgaxis/.style={draw=dgGrey!70, line width=0.6pt, ->},
}

\newenvironment{dgflow}[1][]
  {\begin{tikzpicture}[dgflowbase,#1]}
  {\end{tikzpicture}}

% ---------------------------------------------------------------------
%  Figure wrapper. Used by every figure so the spacing above and below
%  a figure is identical everywhere.
%
%    \dgfigure{label}{short caption}{long caption}{ ...body... }
% ---------------------------------------------------------------------
\newcommand{\dgfigure}[4]{%
  \begin{figure}[htbp]
    \centering
    #4
    \caption[#2]{#3}
    \label{#1}
  \end{figure}}


% ---- Maths -----------------------------------------------------------
\usepackage{amsmath,amssymb}

% ---- Citations and links ---------------------------------------------
% cite sorts and compresses numeric citations: \cite{a,b,c} -> [2]-[4].
% It must load BEFORE hyperref.
\usepackage{cite}
\usepackage[colorlinks=true,linkcolor=black,citecolor=black,
            urlcolor=dgCFT,breaklinks=true]{hyperref}
\usepackage{url}
\urlstyle{same}

% ---- Source-quality markers ------------------------------------------
% Project rule R9: a source that has not been through peer review is
% marked WHERE IT IS USED, not only in the bibliography. A reader
% checking a claim should not have to turn to the back to find out that
% the evidence is a preprint.
% Capitals, not small capitals: Times New Roman has no true small-caps
% face, and neither does the metric-compatible fallback, so \textsc{}
% only produced a font warning and a synthesised shape.
\newcommand{\grey}{\textsuperscript{G}}   % grey literature
\newcommand{\prep}{\textsuperscript{P}}   % preprint, not refereed
% U marks a DIFFERENT thing from G and P. Those two say what kind of
% publication a source is. U says that the thesis could not verify the
% source's own bibliographic data against the article or an index, and
% took it from another work's reference list. A reader who wants to
% follow the citation needs to know that before they try.
\newcommand{\unver}{\textsuperscript{U}}   % bibliographic data unverified

% ---- Front-matter heading -------------------------------------------
% 14 pt bold, centred — the department template's look for the pages
% before Chapter 1. \chapter* would print a three-line monster instead.
%   #1 = as printed   #2 = as it appears in the table of contents
\newcommand{\frontheading}[2]{%
  \phantomsection
  \addcontentsline{toc}{chapter}{#2}%
  \markboth{#2}{#2}%
  \begin{center}\large\bfseries #1\end{center}
  \vspace{0.8em}}

% Blank leaf, as the template has between front-matter sections. It is
% counted and numbered, so both editions and the .docx agree page for page.
\newcommand{\blankleaf}{%
  \clearpage
  \thispagestyle{plain}%
  \null\vfill\begin{center}\small\dgblankleaftext\end{center}\vfill\null
  \clearpage}

% ---- Measurements that do not exist yet ------------------------------
% Project rule R1: never print a number that was not measured. Every
% figure from the author's own benchmark campaign is declared once in
% data/measurements.tex. Until a value is filled in, \meas prints a
% visible marker instead of a plausible-looking number, and
% tools/data-check.py lists what is still outstanding.
\newcommand{\PENDING}{\textbf{\color{dgBFT}[PENDING]}}
\newcommand{\meas}[1]{\csname measvalue@#1\endcsname}
\newcommand{\declaremeasurement}[2]{%
  \expandafter\newcommand\csname measvalue@#1\endcsname{#2}}
\newcommand{\pendingmeasurement}[1]{%
  \expandafter\newcommand\csname measvalue@#1\endcsname{\PENDING}}

\input{style/lang-\probelang}
\input{figures/labels-\probelang}

\pagestyle{empty}
\begin{document}
\thesisnames

\newcommand{\probe}[1]{%
  \begin{center}
    {\ttfamily\small figures/src/#1.tex \quad(\probelang)}\par\vspace{6mm}
    \input{figures/src/#1}
  \end{center}
  \clearpage}

\probe{fig-timeline}
\probe{fig-quorum}
\probe{fig-scaling}
\probe{fig-synchrony}
\probe{fig-equivocation}
\probe{fig-smr}
\end{document}
"
%
%  Both editions must compile with zero errors AND be looked at:
%      pdftoppm -png -r 110 figures/out/probe.pdf /tmp/probe
% =====================================================================
\documentclass[11pt,a4paper]{report}
\providecommand{\probelang}{en}
% =====================================================================
%  common.tex — everything that is the same in both languages.
%
%  Compile with XeLaTeX. pdfLaTeX will not do: the Greek edition needs
%  fontspec, and both editions must use the same font files so the two
%  PDFs look like the same document.
%
%  Language-specific setup lives in style/lang-en.tex and lang-el.tex.
% =====================================================================

% ---- Fonts ----------------------------------------------------------
\usepackage{fontspec}
\usepackage{polyglossia}

% The department requires Times New Roman at 11 pt (Guidelines A.2).
% Times New Roman itself is not redistributable, so the build asks for it
% first and falls back to Liberation Serif, which is metric-compatible:
% same widths, same line breaks, same page count. On a machine that has
% the real font (Word, Overleaf) the first branch wins and nothing moves.
\IfFontExistsTF{Times New Roman}{
  \newcommand{\thesisserif}{Times New Roman}
  \newcommand{\thesissans}{Arial}
  \newcommand{\thesismono}{Courier New}
}{
  \newcommand{\thesisserif}{Liberation Serif}
  \newcommand{\thesissans}{Liberation Sans}
  \newcommand{\thesismono}{Liberation Mono}
}
\setmainfont[Ligatures=TeX]{\thesisserif}
\setsansfont[Ligatures=TeX]{\thesissans}
\setmonofont[Scale=0.90,Ligatures=TeX]{\thesismono}

% ---- Page ------------------------------------------------------------
% 2.54 cm on every side — the margin the department template uses.
\usepackage[a4paper,margin=2.54cm]{geometry}
\usepackage{microtype}
\emergencystretch=1.5em          % prefer loose spacing over an overfull line
\usepackage{setspace}
\onehalfspacing                  % Guidelines A.2: line spacing 1.5
\setlength{\parskip}{0pt}
\setlength{\parindent}{1.2em}

\usepackage{csquotes}
\usepackage{enumitem}
\setlist{itemsep=2pt,parsep=0pt,topsep=4pt}

% ---- Tables and figures ---------------------------------------------
\usepackage{booktabs}
\usepackage{tabularx}
\usepackage{array}
\usepackage{multirow}
\usepackage{graphicx}
\usepackage{xcolor}
\usepackage[font=small,labelfont=bf,textfont=small]{caption}
\usepackage{subcaption}
\usepackage{float}

% ---- Diagrams --------------------------------------------------------
\usepackage{tikz}
\usetikzlibrary{positioning,arrows.meta,calc,fit,backgrounds,shapes.geometric,
                decorations.pathreplacing}
\usepackage{pgfplots}
\pgfplotsset{compat=1.18}
\usepgfplotslibrary{groupplots}
% =====================================================================
%  diagrams.tex — THE diagram standard. One file, every figure.
%
%  The rule this file exists to enforce: a reader should be able to look
%  at any two figures in the thesis and see the same visual language.
%  Same palette, same line weights, same type size, same width.
%
%  There are exactly TWO kinds of figure, and every figure is one of them:
%
%    dgdata  — numbers on axes (pgfplots). Anything measured or computed.
%    dgflow  — structure and sequence (TikZ). Anything explained.
%
%  Nothing in content/ may define a colour, a width, or a font size.
%  tools/diagram-check.py fails the build if it does.
% =====================================================================

% ---------------------------------------------------------------------
%  Palette — Okabe & Ito, chosen because it survives the three common
%  colour-vision deficiencies AND greyscale printing. A thesis gets
%  printed in black and white more often than anyone admits.
%
%  Each colour has ONE meaning across the whole thesis:
% ---------------------------------------------------------------------
\definecolor{dgCFT}{HTML}{0072B2}   % crash-fault protocols            (blue)
\definecolor{dgBFT}{HTML}{D55E00}   % Byzantine-fault protocols     (vermilion)
\definecolor{dgAux}{HTML}{009E73}   % the third series / "correct"     (green)
\definecolor{dgSand}{HTML}{E69F00}  % the fourth series / "attention" (orange)
\definecolor{dgPurple}{HTML}{CC79A7}% the fifth series                (purple)
\definecolor{dgGrey}{HTML}{6E6E6E}  % structure only — never a series   (grey)
\definecolor{dgInk}{HTML}{1A1A1A}   % text and axes
\definecolor{dgPaper}{HTML}{FFFFFF} % label backgrounds over lines

% Colour is never the only difference between two series. Every series
% also gets its own mark and its own dash pattern, so the figure still
% reads on a photocopy. These four cycles stay in lockstep.
\pgfplotscreateplotcyclelist{dgcycle}{%
  {dgCFT,   mark=*,        mark size=1.9pt, solid},
  {dgBFT,   mark=square*,  mark size=1.8pt, dashed},
  {dgAux,   mark=triangle*,mark size=2.1pt, dotted},
  {dgSand,  mark=diamond*, mark size=2.1pt, dashdotted},
  {dgPurple,mark=pentagon*,mark size=1.9pt, densely dashed},
}

% ---------------------------------------------------------------------
%  Fixed geometry. Every figure is the same width, so the eye does not
%  have to re-scale between them.
% ---------------------------------------------------------------------
% Computed at \begin{document}, NOT here: in the preamble \linewidth is
% still the class default (360pt), not the real text block set by
% geometry. Setting it here silently made every figure 10.9cm wide on a
% 15.9cm page.
\newlength{\dgwidth}
\newlength{\dgheight}
\AtBeginDocument{%
  \setlength{\dgwidth}{0.86\linewidth}%
  \setlength{\dgheight}{0.50\dgwidth}%
}

% ---------------------------------------------------------------------
%  KIND 1 — dgdata: numbers on axes
%
%    \begin{dgdata}{xlabel={...}, ylabel={...}}
%      \addplot ... ;
%    \end{dgdata}
%
%  Everything that could drift between figures is set here once.
% ---------------------------------------------------------------------
\pgfplotsset{
  dgbase/.style={
    width=\dgwidth, height=\dgheight,
    cycle list name=dgcycle,
    every axis plot/.append style={line width=0.9pt},
    tick label style={font=\footnotesize, color=dgInk},
    label style={font=\small, color=dgInk},
    title style={font=\small\bfseries, color=dgInk},
    legend style={font=\footnotesize, draw=dgGrey!45, fill=dgPaper,
                  fill opacity=0.92, text opacity=1,
                  cells={anchor=west}, row sep=1pt},
    legend cell align=left,
    axis line style={color=dgGrey!70, line width=0.6pt},
    grid=major,
    major grid style={color=dgGrey!22, line width=0.4pt},
    tick align=outside, tick style={color=dgGrey!70},
    enlarge x limits=0.06, enlarge y limits=0.10,
    axis on top,
  },
  % Bar charts share the same skin but need their own spacing.
  dgbars/.style={
    dgbase,
    ybar, bar width=8pt, enlarge x limits=0.14,
    ymin=0, grid=major, ymajorgrids=true, xmajorgrids=false,
    % "fill=." means "fill with this plot's own cycle colour". Without it
    % the cycle list sets only the line colour and every bar comes out as
    % an empty outline.
    every axis plot/.append style={fill=., fill opacity=0.85,
                                   draw=dgInk!55, line width=0.4pt},
    xtick style={draw=none},
  },
}

\newenvironment{dgdata}[1][]
  {\begin{tikzpicture}\begin{axis}[dgbase,#1]}
  {\end{axis}\end{tikzpicture}}

\newenvironment{dgbarchart}[1][]
  {\begin{tikzpicture}\begin{axis}[dgbars,#1]}
  {\end{axis}\end{tikzpicture}}

% ---------------------------------------------------------------------
%  KIND 2 — dgflow: structure and sequence
%
%    \begin{dgflow}
%      \node[dgbox] (a) {...};
%      \node[dgbox, right=of a] (b) {...};
%      \draw[dgarrow] (a) -- (b) node[dgtag,midway]{...};
%    \end{dgflow}
%
%  Node styles are deliberately few. Five box styles is already too many
%  for a reader to hold; these are four, and each means one thing.
% ---------------------------------------------------------------------
\tikzset{
  dgflowbase/.style={
    node distance=9mm and 13mm,
    every node/.append style={font=\small, text=dgInk, align=center},
    >={Stealth[length=2.4mm,width=1.8mm]},
  },
  % a thing that exists: a node, a component, a layer
  dgbox/.style={
    draw=dgGrey!75, line width=0.6pt, rounded corners=1.5pt,
    fill=dgGrey!7, inner sep=4pt, minimum height=7mm, minimum width=20mm,
  },
  % the same, but this one is the subject of the figure
  dgboxkey/.style={dgbox, draw=dgCFT, line width=0.9pt, fill=dgCFT!8},
  % the same, but this one is where it goes wrong
  dgboxbad/.style={dgbox, draw=dgBFT, line width=0.9pt, fill=dgBFT!8},
  % a moment in time, or a state
  dgstate/.style={
    draw=dgGrey!75, line width=0.6pt, circle, fill=dgPaper,
    inner sep=2.5pt, minimum size=9mm,
  },
  % movement: a message, a step, a dependency
  dgarrow/.style={->, draw=dgGrey!85, line width=0.7pt},
  dgarrowkey/.style={dgarrow, draw=dgCFT, line width=0.9pt},
  dgarrowbad/.style={dgarrow, draw=dgBFT, line width=0.9pt, dashed},
  % a label sitting on a line — the white background is what stops it
  % colliding with the line underneath
  % NOTE: text=, not color=. TikZ's color= sets the fill too, so
  % "color=dgInk" here would override fill=dgPaper and render every tag
  % as dark text on a dark box.
  dgtag/.style={
    font=\footnotesize, text=dgInk, fill=dgPaper, inner sep=1.5pt,
  },
  % an aside: the sentence that explains what the reader is looking at
  dgnote/.style={
    font=\footnotesize\itshape, color=dgGrey, align=left, inner sep=2pt,
  },
  % an axis inside a flow figure (a timeline, a spectrum)
  dgaxis/.style={draw=dgGrey!70, line width=0.6pt, ->},
}

\newenvironment{dgflow}[1][]
  {\begin{tikzpicture}[dgflowbase,#1]}
  {\end{tikzpicture}}

% ---------------------------------------------------------------------
%  Figure wrapper. Used by every figure so the spacing above and below
%  a figure is identical everywhere.
%
%    \dgfigure{label}{short caption}{long caption}{ ...body... }
% ---------------------------------------------------------------------
\newcommand{\dgfigure}[4]{%
  \begin{figure}[htbp]
    \centering
    #4
    \caption[#2]{#3}
    \label{#1}
  \end{figure}}


% ---- Maths -----------------------------------------------------------
\usepackage{amsmath,amssymb}

% ---- Citations and links ---------------------------------------------
% cite sorts and compresses numeric citations: \cite{a,b,c} -> [2]-[4].
% It must load BEFORE hyperref.
\usepackage{cite}
\usepackage[colorlinks=true,linkcolor=black,citecolor=black,
            urlcolor=dgCFT,breaklinks=true]{hyperref}
\usepackage{url}
\urlstyle{same}

% ---- Source-quality markers ------------------------------------------
% Project rule R9: a source that has not been through peer review is
% marked WHERE IT IS USED, not only in the bibliography. A reader
% checking a claim should not have to turn to the back to find out that
% the evidence is a preprint.
% Capitals, not small capitals: Times New Roman has no true small-caps
% face, and neither does the metric-compatible fallback, so \textsc{}
% only produced a font warning and a synthesised shape.
\newcommand{\grey}{\textsuperscript{G}}   % grey literature
\newcommand{\prep}{\textsuperscript{P}}   % preprint, not refereed
% U marks a DIFFERENT thing from G and P. Those two say what kind of
% publication a source is. U says that the thesis could not verify the
% source's own bibliographic data against the article or an index, and
% took it from another work's reference list. A reader who wants to
% follow the citation needs to know that before they try.
\newcommand{\unver}{\textsuperscript{U}}   % bibliographic data unverified

% ---- Front-matter heading -------------------------------------------
% 14 pt bold, centred — the department template's look for the pages
% before Chapter 1. \chapter* would print a three-line monster instead.
%   #1 = as printed   #2 = as it appears in the table of contents
\newcommand{\frontheading}[2]{%
  \phantomsection
  \addcontentsline{toc}{chapter}{#2}%
  \markboth{#2}{#2}%
  \begin{center}\large\bfseries #1\end{center}
  \vspace{0.8em}}

% Blank leaf, as the template has between front-matter sections. It is
% counted and numbered, so both editions and the .docx agree page for page.
\newcommand{\blankleaf}{%
  \clearpage
  \thispagestyle{plain}%
  \null\vfill\begin{center}\small\dgblankleaftext\end{center}\vfill\null
  \clearpage}

% ---- Measurements that do not exist yet ------------------------------
% Project rule R1: never print a number that was not measured. Every
% figure from the author's own benchmark campaign is declared once in
% data/measurements.tex. Until a value is filled in, \meas prints a
% visible marker instead of a plausible-looking number, and
% tools/data-check.py lists what is still outstanding.
\newcommand{\PENDING}{\textbf{\color{dgBFT}[PENDING]}}
\newcommand{\meas}[1]{\csname measvalue@#1\endcsname}
\newcommand{\declaremeasurement}[2]{%
  \expandafter\newcommand\csname measvalue@#1\endcsname{#2}}
\newcommand{\pendingmeasurement}[1]{%
  \expandafter\newcommand\csname measvalue@#1\endcsname{\PENDING}}

\input{style/lang-\probelang}
\input{figures/labels-\probelang}

\pagestyle{empty}
\begin{document}
\thesisnames

\newcommand{\probe}[1]{%
  \begin{center}
    {\ttfamily\small figures/src/#1.tex \quad(\probelang)}\par\vspace{6mm}
    \input{figures/src/#1}
  \end{center}
  \clearpage}

\probe{fig-timeline}
\probe{fig-quorum}
\probe{fig-scaling}
\probe{fig-synchrony}
\probe{fig-equivocation}
\probe{fig-smr}
\end{document}
